\chapter{原子核与放射性}\label{ch:g11:nucleus-radioactivity}

此刻，你体内每一秒都有约八千个\omterm{def:g11:fundamental-interactions:ladder}{原子核}
“爆裂”。无人攻击它们：有些\omterm{def:g11:fundamental-interactions:ladder}{原子核}天生不稳定，每一个
迟早都会自行转变，射出一个碎片，速率达到光速的相当大
一部分。本章将回到 $10^{-15}$ 米尺度的
“底层”（见\cref{ch:g11:fundamental-interactions}），看看哪些\omterm{def:g11:fundamental-interactions:ladder}{原子核}脆弱、如何破裂，以及
它们有规律的“急躁”如何给洞穴定年，又如何为天花板“站岗”。

\section{原子核、核子与同位素}

\begin{definition}[原子核记号]\label{def:g11:nucleus-radioactivity:notation}
\omterm{def:g11:fundamental-interactions:ladder}{原子核}含 $Z$ 个质子（每个带电 $+e$）和 $N$ 个中子（
不带电），共 $A = Z + N$ 个\omterm{def:g11:fundamental-interactions:ladder}{核子}，记为 ${}^{A}_{Z}\mathrm{X}$，
其中 $\mathrm X$ 为化学符号，$Z$ 为\emph{原子序数}\index{原子序数}，
$A$ 为\emph{质量数}\index{质量数}。由 $A$ 和 $Z$ 确定的一类\omterm{def:g11:fundamental-interactions:ladder}{原子核}
称为\emph{核素}\index{核素}；质子数 $Z$ 相同而 $A$ 不同的核素，
互为\emph{同位素}\index{同位素}，
同属一种元素：化学性质相同（电子云只感受到 $Ze$），
而\omterm{def:g11:fundamental-interactions:ladder}{原子核}和稳定性不同。
\end{definition}

\begin{example}[氢与碳]\label{ex:g11:nucleus-radioactivity:hc}
氢有三种\omterm{def:g11:nucleus-radioactivity:notation}{同位素}：${}^{1}_{1}\mathrm{H}$（单个质子），
${}^{2}_{1}\mathrm{H}$（氘：一个质子、一个中子），以及不稳定的
${}^{3}_{1}\mathrm{H}$（氚：两个中子）。天然碳主要是
${}^{12}_{6}\mathrm{C}$（$6$ 个质子、$6$ 个中子），另有 $1.1\%$ 为
${}^{13}_{6}\mathrm{C}$，每 $10^{12}$ 个\omterm{def:g11:fundamental-interactions:ladder}{原子}中有一个不稳定的
${}^{14}_{6}\mathrm{C}$（$8$ 个中子），是\cref{ex:g11:nucleus-radioactivity:dating}的“主角”。
\end{example}

\section{稳定与不稳定的原子核}

\begin{proposition}[稳定谷]\label{prop:g11:nucleus-radioactivity:valley}
稳定\omterm{def:g11:nucleus-radioactivity:notation}{核素}位于 $(Z, N)$ 平面中的一条狭窄带状区域：轻核满足 $N \approx Z$，
随后中子过剩逐渐增大，直至 $N \approx 1.5\,Z$。
凡偏离该带的\omterm{def:g11:nucleus-radioactivity:notation}{核素}——中子过多、质子过多，或
过重（超过铅，$Z = 82$，已无真正稳定的\omterm{def:g11:nucleus-radioactivity:notation}{核素}）——迟早都会转变。
\end{proposition}
\admitted

\begin{remark}
精确确定该带需要量子力学——留给大学阶段的各分册。
这一规律正是\cref{ch:g11:fundamental-interactions}的“账本”：强
“胶水”只束缚近邻，\omterm{def:g11:fundamental-interactions:charge}{库仑}斥力却贯穿整个\omterm{def:g11:fundamental-interactions:ladder}{原子核}，因此
重核需要额外的中子“胶水”；然而中子过剩也不稳定
（这由\omterm{def:g11:fundamental-interactions:weak}{弱相互作用}负责）。
\end{remark}

\begin{omfigure}
\begin{tikzpicture}[font=\small, x=0.048cm, y=0.033cm]
  \draw[-Stealth] (0,0) -- (108,0) node[below left=1pt] {$Z$};
  \draw[-Stealth] (0,0) -- (0,158) node[below right=1pt] {$N$};
  \draw (40,0) -- (40,-4) node[below] {$40$} (0,100) -- (-2,100) node[left] {$100$};
  \draw[dashed] (0,0) -- (105,105) node[right] {$N = Z$};
  \draw[omDef, line width=6pt, opacity=0.4] plot[smooth]
    coordinates {(1,1) (20,22) (40,50) (60,82) (83,126)};
  \draw[omDef, -Stealth] (52,30) node[below] {稳定核素} -- (57,72);
  \node at (40,120) {中子过多：$\beta^-$};
  \node at (82,18) {质子过多：$\beta^+$};
  \node[left] at (83,145) {过重：$\alpha$};
\end{tikzpicture}

{\small 稳定\omterm{def:g11:nucleus-radioactivity:notation}{核素}紧贴 $N = Z$，随后向富中子方向偏移；带的上方
\omterm{def:g11:fundamental-interactions:ladder}{原子核}发生 $\beta^-$ 衰变，带的下方发生 $\beta^+$ 衰变，超过铋之后为 $\alpha$ 衰变。}
\end{omfigure}

\begin{definition}[放射性]\label{def:g11:nucleus-radioactivity:radioactivity}
\emph{放射性}\index{放射性}是指不稳定\omterm{def:g11:fundamental-interactions:ladder}{原子核}
自发转变为另一种\omterm{def:g11:nucleus-radioactivity:notation}{核素}，同时放出辐射。它
是\emph{随机}的：没有任何迹象预示下一个衰变的\omterm{def:g11:fundamental-interactions:ladder}{原子核}，且没有
\omterm{def:g10:pressure:pressure}{压强}、温度或化学作用能够加快或延缓它；只有大量
\omterm{def:g11:fundamental-interactions:ladder}{原子核}组成的总体才可预测（见\cref{def:g11:nucleus-radioactivity:halflife}）。
\end{definition}

\begin{remark}[贝可勒尔与居里夫妇]\label{rem:g11:nucleus-radioactivity:curie}
1896 年，亨利·\omterm{def:g11:nucleus-radioactivity:activity}{贝可勒尔}发现铀盐能使照相底片感光，
隔着黑纸、放在关闭的抽屉里，且无需外界供给\omterm{def:g10:energy-conservation:energy}{能量}。玛丽与
皮埃尔·居里证明这一效应源于\omterm{def:g11:fundamental-interactions:ladder}{原子}本身，测量了它，并分离出钋
和镭，其\omterm{def:g11:nucleus-radioactivity:radioactivity}{放射性}比铀强百万倍。\emph{\omterm{def:g11:nucleus-radioactivity:radioactivity}{放射性}}一词由
玛丽·居里提出，她还获得了两项诺贝尔奖。
\end{remark}

\section{α、β 与 γ 射线}

\begin{definition}[三种历史上的射线]\label{def:g11:nucleus-radioactivity:abg}
三种射线，人们尚未弄清其本质时就已为其命名：
\begin{itemize}
  \item \emph{α 射线}\index{α 射线}（$\alpha$）：氦\omterm{def:g11:fundamental-interactions:ladder}{原子核}
        ${}^{4}_{2}\mathrm{He}$，整体射出——质量大、带两个\omterm{def:g10:orders-of-magnitude:unit}{单位}正电荷，
        一张纸或几厘米空气即可挡住；
  \item \emph{β 射线}\index{β 射线}（$\beta^-$）：一个电子
        ${}^{\;0}_{-1}\mathrm{e}$，在\omterm{def:g11:fundamental-interactions:ladder}{原子核}内产生并射出——
        质量小、速度快，几毫米厚的铝板即可挡住；
  \item \emph{γ 射线}\index{γ 射线}（$\gamma$）：电磁
        辐射，\omterm{def:g10:energy-conservation:energy}{能量}远高于可见光——几乎不会被完全挡住，只能
        被\emph{衰减}\index{衰减}：一厘米厚的铅吸收其一半。
\end{itemize}
\end{definition}

\begin{omfigure}
\begin{tikzpicture}[font=\small]
  \fill[omExo] (-0.15,-0.3) rectangle (0,3.0); \node[left] at (-0.15,1.4) {放射源};
  \draw[omProp, thick] (2.2,-0.3) -- (2.2,3.0) node[above] {纸};
  \draw[omProp, line width=2.5pt] (4.4,-0.3) -- (4.4,3.0) node[above] {铝};
  \fill[omExo] (6.4,-0.3) rectangle (6.9,3.0); \node[above] at (6.65,3.0) {铅};
  \draw[-Stealth, omDef, thick] (0,2.4) -- (2.15,2.4) node[midway, above] {$\alpha$};
  \draw[-Stealth, omDef, thick] (0,1.4) -- (4.35,1.4) node[pos=0.28, above] {$\beta^-$};
  \draw[omThm, thick, decorate,
        decoration={snake, amplitude=0.5mm, segment length=3mm}] (0,0.5) -- (6.4,0.5);
  \draw[-Stealth, omThm, thick, dashed] (6.9,0.5) -- (7.9,0.5);
  \node[omThm, above] at (1.1,0.62) {$\gamma$};
\end{tikzpicture}

{\small $\alpha$ 在纸中即被挡住，$\beta^-$ 则需几毫米厚的铝板；
$\gamma$ 即便是铅，也只能使其\omterm{def:g11:nucleus-radioactivity:abg}{衰减}。}
\end{omfigure}

\begin{proposition}[核反应方程中的守恒]\label{prop:g11:nucleus-radioactivity:soddy}
在任一核转变中，总\omterm{def:g11:nucleus-radioactivity:notation}{质量数} $A$ 与总
电荷——即各下标之和，发射出的
电子计为 $-1$，在转变前后保持不变。
\end{proposition}
\admitted

\begin{remark}
电荷守恒已在\cref{ch:g11:fundamental-interactions}中排除了这种可能；
$A$ 守恒表明\omterm{def:g11:fundamental-interactions:ladder}{核子}只会相互转化，
既不会凭空产生，也不会消失——这是留给大学分册的严谨记账。
\end{remark}

\begin{method}[配平衰变方程]\label{met:g11:nucleus-radioactivity:balance}
\begin{enumerate}
  \item 每个粒子都要同时标出上下标：$\alpha = {}^{4}_{2}\mathrm{He}$，
        $\beta^- = {}^{\;0}_{-1}\mathrm{e}$，$\beta^+ = {}^{0}_{+1}\mathrm{e}$；$\gamma$ 满足 $A = Z = 0$。
  \item 令各 $A$ 之和相等，再令各 $Z$ 之和相等
        （见\cref{prop:g11:nucleus-radioactivity:soddy}）；解方程，再由
        子核的 $Z$ 在元素\omterm{def:g10:signals-and-waves:period}{周期}表中查得其元素。
\end{enumerate}
\end{method}

\begin{example}[α 衰变]\label{ex:g11:nucleus-radioactivity:alpha}
镭-226（即居里夫妇所研究的镭）是 $\alpha$ 发射体：
${}^{226}_{88}\mathrm{Ra} \to {}^{A}_{Z}\mathrm{Y} + {}^{4}_{2}\mathrm{He}$
由此得 $A = 226 - 4 = 222$，$Z = 88 - 2 = 86$；元素 $86$ 为氡，故
${}^{226}_{88}\mathrm{Ra} \to {}^{222}_{86}\mathrm{Rn} + {}^{4}_{2}\mathrm{He}$
——通风不良的地窖中的\omterm{def:g11:nucleus-radioactivity:radioactivity}{放射性}气体。
\end{example}

\begin{example}[负 β 衰变]\label{ex:g11:nucleus-radioactivity:beta}
碳-14 位于稳定带上方（中子数 $8$，质子数 $6$）：
${}^{14}_{6}\mathrm{C} \to {}^{14}_{7}\mathrm{N} + {}^{\;0}_{-1}\mathrm{e}$。
检验：$14 = 14 + 0$，$6 = 7 - 1$；一个中子变成了一个质子加
射出的电子，这便是\omterm{def:g11:fundamental-interactions:weak}{弱相互作用}在起作用。
\end{example}

\begin{remark}[正 β 衰变与 γ 辐射]\label{rem:g11:nucleus-radioactivity:betaplus}
在稳定带下方，镜像衰变 $\beta^+$ 使一个质子转变为中子
加上一个\emph{正电子}\index{正电子} ${}^{0}_{+1}\mathrm{e}$，即电子的反粒子：
${}^{18}_{9}\mathrm{F} \to {}^{18}_{8}\mathrm{O} + {}^{0}_{+1}\mathrm{e}$，
这种衰变是 PET 扫描仪的主力。多数衰变之后，子核刚生成时处于激发态，
便通过发射一个 $\gamma$ 辐射而趋于稳定。
\end{remark}

\begin{omfigure}
\begin{tikzpicture}[font=\small]
  \node (U)  at (0,0)   {${}^{238}_{92}\mathrm{U}$};
  \node (Th) at (2.6,0) {${}^{234}_{90}\mathrm{Th}$};
  \node (Pa) at (5.2,0) {${}^{234}_{91}\mathrm{Pa}$}; \node (U4) at (7.8,0) {${}^{234}_{92}\mathrm{U}$};
  \draw[-Stealth, omDef, thick] (U) -- (Th) node[midway, above] {$\alpha$};
  \draw[-Stealth, omProp, thick] (Th) -- (Pa) node[midway, above] {$\beta^-$};
  \draw[-Stealth, omProp, thick] (Pa) -- (U4) node[midway, above] {$\beta^-$};
\end{tikzpicture}

{\small 铀-238 衰变链的前几步；另有十一次衰变
（$\alpha$ 与 $\beta^-$）最终止于稳定的 ${}^{206}_{82}\mathrm{Pb}$。}
\end{omfigure}

\section{半衰期与放射性活度}

\begin{definition}[半衰期]\label{def:g11:nucleus-radioactivity:halflife}
某\omterm{def:g11:nucleus-radioactivity:notation}{核素}的\emph{半衰期}\index{半衰期} $T_{1/2}$ 是指
经过此时间后，任一大量样本中有一半已发生衰变。每再经过一个半衰期，
剩余部分每次减半：$N_0$ 个\omterm{def:g11:fundamental-interactions:ladder}{原子核}中，$N_0/2$ 个在 $T_{1/2}$ 时剩下，$N_0/4$
个在 $2\,T_{1/2}$ 时，$N_0/8$ 个在 $3\,T_{1/2}$ 时；一般地，$N_0/2^{\,n}$ 个在经过 $n$ 个半衰期后剩下。
\end{definition}

\begin{remark}
单个\omterm{def:g11:fundamental-interactions:ladder}{原子核}的衰变具有随机性，而 $10^{23}$ 个\omterm{def:g11:fundamental-interactions:ladder}{原子核}则如钟表般精确：好比一大群
抛硬币者，逐一难以预测，每一轮都淘汰一半。
若要在整数个\omterm{def:g11:nucleus-radioactivity:halflife}{半衰期}\emph{之间}插值，便要用到
\omterm{def:g10:orders-of-magnitude:scinot}{指数}函数——留待明年，在数学分册中学习。
\end{remark}

\begin{example}[二十个数量级]\label{ex:g11:nucleus-radioactivity:range}
\begin{center}\small\begin{tabular}{llll}
钋-214 & $\alpha$ & \qty{1.6e-4}{s} & 镭衰变链中的一环 \\
锝-99m & $\gamma$ & $6.0$ 小时 & 医学成像 \\
镅-241 & $\alpha$ & $432$ 年 & 烟雾探测器 \\
碳-14 & $\beta^-$ & \num{5700} 年 & 考古学 \\
钾-40 & $\beta$ & \num{1.25e9} 年 & 每根香蕉中都有 \\
铀-238 & $\alpha$ & \num{4.5e9} 年 & 地球内部热量 \\
\end{tabular}\end{center}
\end{example}

\begin{example}[碳-14 测年]\label{ex:g11:nucleus-radioactivity:dating}
活体与大气不断交换碳，因此其中所含
碳-14 的比例与大气中的相同；死亡后摄入停止，该比例每
隔 \num{5700} 年减半。取自一处有壁画洞穴的木炭，
其比例仅为活体比例的四分之一：四分之一即二分之一的二分之一，
故此火烧于两个\omterm{def:g11:nucleus-radioactivity:halflife}{半衰期}之前，即 $2 \times \num{5700} = \num{11400}$ 年前。
\end{example}

\begin{definition}[放射性活度]\label{def:g11:nucleus-radioactivity:activity}
样品的\emph{放射性活度}\index{放射性活度} $\mathcal A$ 是指其
每秒衰变的次数，以\emph{贝可勒尔}\index{贝可勒尔}为\omterm{def:g10:orders-of-magnitude:unit}{单位}：
$\qty{1}{Bq} = 1$ 即每秒一次衰变。尚存的\omterm{def:g11:fundamental-interactions:ladder}{原子核}越少，
衰变便越少，故活度也在每个\omterm{def:g11:nucleus-radioactivity:halflife}{半衰期}减半。
\end{definition}

\begin{omfigure}
\begin{tikzpicture}[font=\small, x=1.1cm, y=3.0cm]
  \draw[-Stealth] (0,0) -- (5.7,0) node[below left=1pt] {$t$（以 $T_{1/2}$ 计）};
  \draw[-Stealth] (0,0) -- (0,1.15) node[right] {$\mathcal A/\mathcal A_0$};
  \foreach \n/\v in {0/1, 1/0.5, 2/0.25, 3/0.125, 4/0.0625}{
    \filldraw[fill=omDef!35, draw=omDef] (\n+0.15,0) rectangle (\n+0.85,\v);
    \node[below] at (\n+0.5,0) {$\n$};}
  \node[left] at (0,1) {$1$}; \node[left] at (0,0.5) {$\tfrac12$};
  \node[left] at (0,0.25) {$\tfrac14$};
  \draw[dashed, omExo] (0,0.5) -- (1.15,0.5) (0,0.25) -- (2.15,0.25);
\end{tikzpicture}

{\small 递减阶梯：经过 $n$ 个\omterm{def:g11:nucleus-radioactivity:halflife}{半衰期}后，活度为
$\mathcal A_0/2^{\,n}$ ——三个\omterm{def:g11:nucleus-radioactivity:halflife}{半衰期}后为八分之一，十个\omterm{def:g11:nucleus-radioactivity:halflife}{半衰期}后约为千分之一。}
\end{omfigure}

\begin{example}[数量级]\label{ex:g11:nucleus-radioactivity:orders}
一切事物都带有轻微的\omterm{def:g11:nucleus-radioactivity:radioactivity}{放射性}。香蕉：约 \qty{15}{Bq}
（钾-40）。人体：约 \qty{8000}{Bq}——本章的引子，
相当于约五百根香蕉的\omterm{def:g11:nucleus-radioactivity:radioactivity}{放射性}。一立方米花岗岩：约
\qty{3e6}{Bq}（来自铀及其衰变链，地窖中的氡即由此而来）。一次锝
注射：约 \qty{5e8}{Bq}，数日内即消失殆尽。
\end{example}

\section{服务与危险}

\begin{example}[不稳定原子核的三种用途]\label{ex:g11:nucleus-radioactivity:uses}
\emph{医学}：锝-99m 是纯 $\gamma$ 发射体，
$T_{1/2} = 6$ 小时，将其附着于靶器官会吸收的\omterm{def:g11:fundamental-interactions:ladder}{分子}上；相机
从体外拍摄 $\gamma$ 射线，次日示踪剂已
基本消失。放射治疗则反其道而行之，聚焦射束
以摧毁肿瘤。\emph{断代}：碳-14 可用于测定木材、骨骼和布料的年代，
可追溯约 \num{50000} 年；铀-238 用于岩石——乃至地球，
其年龄约为 \num{4.5e9} 年。\emph{烟雾报警器}：一小粒镅-241 使小室内的空气电离，
产生微弱电流；烟雾阻断电流，
报警器随即响起——$\alpha$ 粒子在几厘米空气中即被阻挡。
\end{example}

\begin{remark}[剂量与希沃特]\label{rem:g11:nucleus-radioactivity:dose}
辐射会将电子从\omterm{def:g11:fundamental-interactions:ladder}{分子}中剥离，并可损伤 DNA。生物
危害以\emph{剂量}\index{剂量}衡量，\omterm{def:g10:orders-of-magnitude:unit}{单位}为
\emph{希沃特}\index{希沃特}（\unit{Sv}），它按
每\omterm{def:g10:orders-of-magnitude:unit}{千克}组织吸收的\omterm{def:g10:energy-conservation:energy}{能量}，并按各类辐射的危害程度加权。
有三条规则：距离、屏蔽（$\alpha$：皮肤即可；
除非放射源被\emph{吸入}，这正是氡的危害所在；$\beta$：铝；
$\gamma$：铅）以及时间。\omterm{rem:g11:nucleus-radioactivity:dose}{剂量}学本身属于大学课程内容。
\end{remark}

\section{练习}

\begin{exercise}[$\star$]\label{exo:g11:nucleus-radioactivity:1}
求下列\omterm{def:g11:fundamental-interactions:ladder}{原子核}的质子数与中子数：${}^{16}_{8}\mathrm{O}$、
${}^{27}_{13}\mathrm{Al}$、${}^{60}_{27}\mathrm{Co}$、${}^{235}_{92}\mathrm{U}$
和 ${}^{3}_{1}\mathrm{H}$。
\end{exercise}

\begin{exercise}[$\star$]\label{exo:g11:nucleus-radioactivity:2}
在 ${}^{14}_{6}\mathrm{C}$、${}^{14}_{7}\mathrm{N}$、${}^{12}_{6}\mathrm{C}$、
${}^{13}_{6}\mathrm{C}$ 和 ${}^{16}_{8}\mathrm{O}$ 中，哪些互为
\omterm{def:g11:nucleus-radioactivity:notation}{同位素}？为何 ${}^{14}_{6}\mathrm{C}$ 与 ${}^{14}_{7}\mathrm{N}$ \emph{不是}
\omterm{def:g11:nucleus-radioactivity:notation}{同位素}，尽管二者\omterm{def:g11:nucleus-radioactivity:notation}{质量数}相同？
\end{exercise}

\begin{exercise}[$\star$]\label{exo:g11:nucleus-radioactivity:3}
钋-210，${}^{210}_{84}\mathrm{Po}$，是 $\alpha$ 发射体。写出
衰变方程，并说明子核的名称（$Z = 81$：铊，$82$：铅，$83$：
铋）。
\end{exercise}

\begin{exercise}[$\star$]\label{exo:g11:nucleus-radioactivity:4}
写出 $\beta^-$ 衰变方程：钴-60（${}^{60}_{27}\mathrm{Co}$）与
碘-131（${}^{131}_{53}\mathrm{I}$）；子核为：$Z = 28$ 镍，
$Z = 54$ 氙。
\end{exercise}

\begin{exercise}[$\star$]\label{exo:g11:nucleus-radioactivity:5}
某碘-131 源（$T_{1/2} = 8.0$ 天）的活度为 \qty{800}{MBq}，
今日。$16$ 天后其活度为多少？$32$ 天后呢？
\end{exercise}

\begin{exercise}[$\star\star$]\label{exo:g11:nucleus-radioactivity:6}
四种辐射穿过强\omterm{def:g11:magnetic-fields:bfield}{磁场}，再穿过屏蔽物。判别每一种：
(a) 被纸挡住，在\omterm{def:g11:magnetic-fields:bfield}{磁场}中几乎不偏转；(b) 能穿过纸，
被 \qty{4}{mm} 铝挡住，偏转明显；(c) 两者均能穿过，
仅被铅\omterm{def:g11:nucleus-radioactivity:abg}{衰减}，不偏转；(d) 与 (b) 相似，但偏转方向相反。
\end{exercise}

\begin{exercise}[$\star\star$]\label{exo:g11:nucleus-radioactivity:7}
氡-222（${}^{222}_{86}\mathrm{Rn}$）开启一条快速衰变链：先是一次 $\alpha$
衰变，接着又是一次 $\alpha$ 衰变，然后是一次 $\beta^-$ 衰变。写出三个方程
（$Z = 82$：铅，$83$：铋，$84$：钋）。
\end{exercise}

\begin{exercise}[$\star\star$]\label{exo:g11:nucleus-radioactivity:8}
氟-18（${}^{18}_{9}\mathrm{F}$）是 PET 扫描仪所用的示踪剂；稳定的
氟为 ${}^{19}_{9}\mathrm{F}$。它位于稳定带的
哪一侧？预测其衰变方式，并写出方程（$Z = 8$：氧）。
\end{exercise}

\begin{exercise}[$\star\star$]\label{exo:g11:nucleus-radioactivity:9}
一名患者接受 \qty{500}{MBq} 的锝-99m（$T_{1/2} = 6.0$ 小时），
于 08:00 注射。次日早上 08:00 的活度为多少？经过多少
小时后，活度首次降至 \qty{1}{MBq} 以下？
\end{exercise}

\begin{exercise}[$\star\star$]\label{exo:g11:nucleus-radioactivity:10}
一块取自泥炭沼泽的骨骼，其碳-14 的比例为活骨骼的
八分之一（$T_{1/2} = \num{5700}$ 年）。该骨骼有多少岁？又为何
碳-14 无法用于测定恐龙骨骼的年代（其年龄约为 $10^{8}$ 年）？
\end{exercise}

\begin{exercise}[$\star\star$]\label{exo:g11:nucleus-radioactivity:11}
利用稳定带，预测衰变方式（$\alpha$、$\beta^-$ 或
$\beta^+$）。${}^{14}_{6}\mathrm{C}$、${}^{11}_{6}\mathrm{C}$ 与
${}^{238}_{92}\mathrm{U}$ 的衰变方式，每项用一行说明理由（稳定的碳：
${}^{12}_{6}\mathrm{C}$、${}^{13}_{6}\mathrm{C}$）。
\end{exercise}

\begin{exercise}[$\star\star\star$]\label{exo:g11:nucleus-radioactivity:12}
铀-238 衰变链经过多步后，终止于稳定的
${}^{206}_{82}\mathrm{Pb}$；每一步都是 $\alpha$ 衰变或 $\beta^-$ 衰变。仅利用
$A$ 守恒，求 $\alpha$ 衰变的步数；再利用
$Z$ 守恒，求 $\beta^-$ 衰变的步数。
\end{exercise}

\begin{exercise}[$\star\star\star$]\label{exo:g11:nucleus-radioactivity:13}
某医院的钴-60放射源（$T_{1/2} = 5.3$ 年）的\omterm{def:g11:nucleus-radioactivity:activity}{放射性活度}为
\qty{6.4}{TBq}；规定\omterm{def:g11:nucleus-radioactivity:activity}{放射性活度}低于 \qty{0.1}{TBq} 时方可处置。经过多少
年后才可处置？
\end{exercise}

\begin{exercise}[$\star\star\star$]\label{exo:g11:nucleus-radioactivity:14}
你的身体的\omterm{def:g11:nucleus-radioactivity:activity}{放射性活度}约为 \qty{8000}{Bq}。每天有多少个\omterm{def:g11:fundamental-interactions:ladder}{原子核}
发生衰变？在 $80$ 年的一生中共有多少？一根香蕉在你消化它时约增加 \qty{15}{Bq}，
请用一句话评论那些让人对每一个\omterm{def:g11:nucleus-radioactivity:activity}{贝可勒尔}都感到恐惧的新闻标题。
\end{exercise}

\begin{exercise}[$\star\star\star$]\label{exo:g11:nucleus-radioactivity:15}
一个烟雾探测器内含 \qty{0.2}{\micro\gram} 的镅-241
（$T_{1/2} = 432$ 年），\omterm{def:g11:nucleus-radioactivity:activity}{放射性活度}约为 \qty{25}{kBq}。每天有多少次
衰变？使用十年后，镅已衰变的比例是
远少于一半、约一半，还是远多于一半——放射源耗尽
是否正是探测器需要更换的原因？最后，为何这个 $\alpha$ 源
装在天花板上是无害的，但若外壳被磨碎、粉尘被吸入，则十分危险？
\end{exercise}

\section{问题：木乃伊、反应堆与烟雾探测器}

\begin{problem}[{周末问题——一个不稳定原子核的三种职业：
用半衰期给木乃伊定年、配平铀的衰变家谱、核查烟雾探测器，
以及恒星为何必定会消亡}]
\label{pb:g11:nucleus-radioactivity:1}
三种不稳定\omterm{def:g11:fundamental-interactions:ladder}{原子核}，三项工作：碳-14 为每一株死去的植物
和动物打上时间标记；铀-238 从内部加热行星，镅-241 监测
天花板。数据：$T_{1/2} = \num{5700}$ 年（C-14），\num{4.5e9} 年
（U-238），$432$ 年（Am-241）；活体碳每克每分钟显示 $13.6$ 次衰变；
\omterm{def:g11:nucleus-radioactivity:notation}{原子序数}：N $7$、O $8$、Pb $82$、Bi $83$、Po $84$、Rn $86$、
Ra $88$、Th $90$、Pa $91$、U $92$、Np $93$、Am $95$。

\medskip
\noindent\textbf{第一部分——木乃伊。}
\begin{enumerate}
  \item 写出 ${}^{12}_{6}\mathrm{C}$ 与
        ${}^{14}_{6}\mathrm{C}$ 的组成（质子数、中子数）。为何它们在体内的化学性质相同？
  \item 为什么活体内碳-14 的比例保持不变，
        而死亡后这一比例为何开始下降？
  \item 写出碳-14 的衰变方程。其射线能否穿透
        木乃伊的亚麻裹布？
  \item 取自亚麻裹布的一克碳每分钟显示 $6.8$ 次衰变。
        这是活体速率的几分之几？请推算木乃伊的年代。
  \item 私人收藏的珍品，第二具“木乃伊”，
        每克每分钟显示 $13.2$ 次衰变。结论如何？
  \item 估算一个经历了十个\omterm{def:g11:nucleus-radioactivity:halflife}{半衰期}的样品的衰变速率，并说明
        为何碳定年法在大约 \num{57000} 年之后便逐渐失效。
\end{enumerate}

\medskip
\noindent\textbf{第二部分——草地之下的反应堆。}
\begin{enumerate}[resume]
  \item 写出 ${}^{238}_{92}\mathrm{U}$ 的组成。为什么重\omterm{def:g11:fundamental-interactions:ladder}{原子核}
        需要如此多的中子富余？
  \item 写出它的 $\alpha$ 衰变方程。
  \item 其子核随后发生 $\beta^-$ 衰变，孙核
        再次发生 $\beta^-$ 衰变。写出两个方程，哪种元素再次出现？
  \item 衰变链终止于 ${}^{206}_{82}\mathrm{Pb}$。数出其中的 $\alpha$
        衰变步数，再数出 $\beta^-$ 衰变步数。
  \item 地球形成于 \num{4.5e9} 年前。其原始铀-238 中
        今天仍有几分之几残存？再过 $5$ 十亿年后呢？
  \item 用一句话回答：这种埋藏于地下的衰变热在地表
        驱动了什么？（想一想火山与漂移的大陆。）
\end{enumerate}

\medskip
\noindent\textbf{第三部分——烟雾探测器。}
\begin{enumerate}[resume]
  \item 写出 $\alpha$ 衰变方程，其母核为 ${}^{241}_{95}\mathrm{Am}$。
  \item 其\omterm{def:g11:nucleus-radioactivity:activity}{放射性活度}为 \qty{33}{kBq}：每秒有多少次衰变？每天呢？
  \item 解释探测器的原理：$\alpha$ 粒子使探测室内的空气
        发生了什么变化？烟雾又改变了什么？
  \item 为何此处选用 $\alpha$ 放射源是合适的，而
        \omterm{def:g11:nucleus-radioactivity:activity}{放射性活度}相同的 $\gamma$ 放射源却是最差的选择？
  \item 经过一个\omterm{def:g11:nucleus-radioactivity:halflife}{半衰期}（$432$ 年！）后，\omterm{def:g11:nucleus-radioactivity:activity}{放射性活度}将为
        \qty{16.5}{kBq}；而说明书要求十年后更换。
        放射源是薄弱环节吗？
\end{enumerate}

\medskip
\noindent\textbf{第四部分——为何恒星必定会消亡。}
\begin{enumerate}[resume]
  \item 垂死的恒星锻造并散布了每一种比氦更重的
        元素。\num{9e9} 年前锻造的铀-238 中，有几分之几
        留存至今？为何地球上残存的丰度表明地球
        远比最古老的恒星年轻？
  \item 地球内部的\omterm{def:g11:nucleus-radioactivity:radioactivity}{放射性}热驱动火山活动与板块
        构造运动，并使地核持续翻涌，从而产生
        在\cref{ch:g11:magnetic-fields}中讨论的磁屏障。若一颗行星只含稳定的
        \omterm{def:g11:fundamental-interactions:ladder}{原子}，它会是什么样子？
  \item 用一句话，既是诗也是账目：你身体里的\omterm{def:g11:fundamental-interactions:ladder}{原子}
        锻造于何处？脚下的大地仍由什么来加热？
\end{enumerate}
\end{problem}
